\chapter{Related Work}
The automation of recruitment and labor market analysis has witnessed a paradigm shift, evolving from manual curation to data-driven intelligence, fueled by advances in Natural Language Processing (NLP) and machine learning \cite{tamburri2020data, faliagka2012online, lacic2020recommender}. This chapter provides a critical review of the existing literature and technological approaches that form the foundation of this thesis. We dissect the state-of-the-art in three interconnected areas: (1) Job Posting Analytics (JPA) and skill extraction, (2) the alignment of extracted information to formal ontologies and taxonomies, and (3) the architectural trade-offs between monolithic Large Language Model (LLM)-based pipelines and the modular, lightweight approach advocated in this work. By synthesizing these domains, we identify a crucial gap in the literature: the need for a system that achieves high semantic accuracy and recall while maintaining the low latency, cost-efficiency, and data sovereignty required for a sustainable, production-grade internship platform.

\section{Job Posting Analytics and Skill Extraction}
The extraction of structured, actionable information from the inherently noisy and semi-structured text of job advertisements is a well-established research area, often referred to as Job Posting Analytics (JPA) \cite{tamburri2020data, boselli2018classifying, kurekova2015using}. The core challenges of JPA—identifying and normalizing job titles, extracting required skills, and classifying experience levels—are microcosms of broader Information Extraction (IE) problems, where the goal is to populate a relational database from unstructured text \cite{sarawagi2008information}.

Early approaches to JPA were dominated by knowledge engineering. Rule-based systems, powered by manually curated gazetteers (lists of known entities like ``Java,'' ``Project Management,'' or ``Berlin''), were the primary method for identifying key entities \cite{sarawagi2008information}. These systems offered high precision on the terms they knew and were highly efficient and explainable. However, their fragility became apparent as they fatally failed to capture the high variability, synonymy, and neologisms of natural language in job descriptions. A gazetteer containing ``C++'' would miss a posting requiring ``experience with the C++ programming language,'' and it would be entirely blind to an emergent skill like ``Prompt Engineering'' without a manual update \cite{li2020survey}.

This brittleness catalyzed the adoption of statistical and deep learning methods, which formulate NER and skill extraction as sequence labeling tasks \cite{li2020survey, wadden2019entity}. These models learn to classify tokens like ``React'' as a Skill based on their bidirectional context. A landmark advancement came with the fine-tuning of Transformer-based models, particularly BERT (Bidirectional Encoder Representations from Transformers), for the recruitment domain \cite{bhola2020retrieving, pappas2022jobbert}. General-purpose BERT models, trained on corpora like Wikipedia, understand language but not the intricate semantics of the job market. Domain-specific pre-training, as demonstrated by JobBERT and similar models, significantly improves the accuracy of skill extraction by learning that ``AWS'' is functionally equivalent to ``Amazon Web Services'' and that a term like ``Spark'' can be disambiguated between a big data framework and a creative skill based on context \cite{pappas2022jobbert}. This thesis directly builds upon this finding, advocating for a lightweight, domain-adapted neural model (Layer B) to handle the semantic complexity that defeats a purely rule-based system.

Beyond simple term extraction, a richer line of research explores the latent semantics of skills. Work on skill embeddings, such as Job2Vec, leverages representation learning to map skills into a high-dimensional vector space where functionally related technologies cluster together, capturing labor market dynamics without explicit rule programming \cite{zhao2018job2vec, giabelli2021skills}. Such systems learn that ``React,'' ``Vue,'' and ``Angular'' are semantically related front-end frameworks. A further layer of sophistication involves analyzing the linguistic context of a skill mention. Qin et al. (2018) proposed methods to automatically classify a skill as ``required,'' ``beneficial,'' or ``preferred,'' providing a nuanced, multi-dimensional profile of employer demands rather than a flat list of keywords \cite{qin2018enhancing}. This granular, contextual understanding is the gold standard for job description analysis, and while our current Lexicon-Based Matching (Layer A) provides a flat skill list, the architecture is designed to incorporate such context-aware weighting in its future iterations.

\section{Ontology/Taxonomy Alignment in HR/Recruitment}
A fundamental obstacle in bridging the labor market's ``vocabulary gap'' is the non-trivial task of mapping raw, heterogeneous text to a canonical reference. Standardized taxonomies like ESCO (European Skills, Competences, Qualifications and Occupations) and O*NET provide this common language, a Rosetta Stone for the labor market \cite{peterson2001occupational}. However, the process of Entity Linking or Taxonomy Alignment—mapping a colloquial title like ``Code Ninja'' or a specific term like ``PyTorch'' to a formal node in a hierarchy—remains a persistent research challenge \cite{zhang2022survey}.

Modern approaches to this problem often leverage Recruitment Knowledge Graphs (KGs) to maintain a web of complex, dynamic relationships between entities like job titles, skills, industries, and educational qualifications \cite{shi2020recruitment}. For instance, a KG can explicitly model the ternary relationship: ``Python \textit{is\_a\_prerequisite\_for} Data Scientist (ESCO Occupation).'' These graph-based systems provide rich, inferable context but often incur high maintenance costs. The literature discusses an inherent trade-off between the deep expressiveness of full ontologies and the practical utility of flatter, more navigable taxonomies \cite{staab2009handbook}. Full ontologies can model complex logical relationships, but they suffer from ``ontology drift'': as new technologies like ``LangChain'' or ``Rust'' emerge, the ontology requires sophisticated, often semi-automated matching and evolution strategies to avoid becoming stale \cite{shvaiko2013ontology}.

Consequently, lightweight semantic approaches, which prioritize operational feasibility, have gained traction. These methods often focus on alignment at the embedding level, converting textual mentions and canonical concepts into the same dense vector space and mapping a raw string to its nearest canonical neighbor via cosine similarity \cite{javed2021review, giabelli2021skills}. This is the philosophy adopted in this thesis. Instead of a high-maintenance KG, we use an SBERT-based bi-encoder to map an idiosyncratic title like ``Software Engineering Intern'' to a region in the vector space near the ESCO concept for ``Software Developer.'' This approach directly addresses the specific challenge of normalizing compound and non-standard job titles, a problem explored by Sayfullina et al. (2018), who demonstrated that learning the decomposition of compound titles is key to robust canonicalization \cite{sayfullina2018learning}. Our approach implicitly performs a similar decomposition at the semantic level.

\section{LLM-based Pipelines vs Lightweight Pipelines}
The emergence of Large Language Models (LLMs) has introduced a seismic shift in the IE paradigm, fundamentally altering the ``make or buy'' decision for semantic capabilities. The prevailing methodology has moved from ``pre-train and fine-tune'' (train a specialist model) to ``pre-train, prompt, and predict'' (engineer a prompt for a generalist model) \cite{brown2020language, liu2023pre}. LLMs like GPT-4 exhibit remarkable zero-shot and few-shot capabilities. With advanced prompting techniques, such as Chain-of-Thought (CoT), an LLM can decompose a complex extraction task—for instance, ``identify and categorize the required versus preferred skills from this job description and map the job title to an ESCO level 4 occupation''—and significantly outperform fine-tuned specialist models that were the previous state-of-the-art, all without a single labeled training example for that specific task \cite{wei2022chain}. This represents an astonishingly powerful and flexible tool.

However, a growing body of critical, comparative studies has soberly highlighted the significant drawbacks of an LLM-only architecture for a production-scale pipeline, particularly one governed by the ``Green AI'' ethos \cite{schwartz2020green}. This thesis argues that these drawbacks are not minor engineering obstacles but fundamental limitations that define the appropriate, bounded role of an LLM. The key challenges are:

\begin{itemize}
    \item \textbf{Latency and Throughput:} Processing thousands of fresh scrapes per hour with cloud-based LLM APIs can result in prohibitive latency (seconds per request) and financial cost, directly clashing with the sub-50ms retrieval goal of an interactive user-facing search engine \cite{hsieh2023distilling, menghani2023efficient}. The cost is thus not only a budgetary line item but a bottleneck to scalability.
    \item \textbf{Consistency and Hallucination:} The well-documented tendency of LLMs to ``hallucinate''---generating plausible-sounding but factually incorrect information---is especially pernicious in extraction. A hallucinated ``Kubernetes'' skill or a misformatted output field can corrupt a database and degrade user trust, a known limitation of current generative architectures \cite{ji2023survey, bender2021dangers}. Deterministic rule-based pipelines, in contrast, are perfectly consistent and debuggable.
    \item \textbf{Privacy and Data Sovereignty:} Sending sensitive recruitment data, potentially containing proprietary employer information or personally identifiable data, to external third-party APIs creates a significant legal and ethical risk vector, especially concerning regional data protection regulations like the GDPR \cite{hunkenschroer2022ethics}. A lightweight, local processing pipeline inherently mitigates this, keeping data within the system's sovereign perimeter.
\end{itemize}

The alternative, and the methodological heart of this thesis, is a lightweight hybrid pipeline. This approach combines the perfect precision of heuristic rules for high-confidence, highly structured data (e.g., regex for salary parsing) with optimized small-scale transformer models like DistilBERT for more complex semantic tasks (e.g., title normalization). This strategy is deeply connected to the principle of Knowledge Distillation, where a compact ``student'' model is trained to replicate the performance of a computationally expensive ``teacher'' LLM on a specific task \cite{hsieh2023distilling, menghani2023efficient}. This thesis extends this concept by proposing a \textit{runtime distillation architecture}: an LLM can serve as an offline teacher for high-quality data annotation and initial schema generation, but the production system utilizes a fast, local student model. This is not just a technical choice but a strategic one, adhering to the principles of ``Green AI'' by prioritizing computational parsimony and operational efficiency without sacrificing the semantic depth required to solve the core ``vocabulary gap'' problem \cite{schwartz2020green}. Our contribution is a concrete system design and empirical evaluation proving that such a hybrid, layered approach constitutes a Pareto-optimal solution, balancing the critical trade-off between accuracy, cost, and control for the domain of internship analytics.